% % Reemplace por el título de su capítulo.
\chapter{El camino}

% % Reemplace por un epígrafe genérico y su atribución.
\epigraph{Caminar despacio también es avanzar.}{Atribuido a un refrán popular}

% % Reemplace por la escena inicial: dónde empieza el capítulo.
La mañana empieza con una taza caliente sobre la mesa y una lista corta de
tareas que cabe en una tarjeta. El primero en levantarse abre la ventana,
mira la calle vacía y decide que hoy basta con hacer bien tres cosas. Como
dice la casa, \enquote{lo poco, hecho con cuidado, rinde el doble}.

% % Reemplace por el desarrollo: qué cambia durante el día.
A media mañana el plan se tuerce un poco: una visita breve, una llamada
larga y un recado que nadie esperaba. En lugar de correr, el cuaderno
recibe una línea nueva y la tarde se ordena otra vez. El truco está en
mover solo lo necesario y dejar lo demás para mañana sin culpa.

\begin{center}
  * * *
\end{center}

% % Reemplace por el cierre: qué queda aprendido al final del día.
Por la noche la lista muestra dos marcas hechas y una pendiente. No es una
victoria grande, pero es una medida honesta del día. La última línea del
cuaderno dice que mañana el camino sigue, con los mismos zapatos y con un
poco más de calma.
